Modern high-energy physics and gravitational-wave astronomy, despite being very different in their experimental settings, share a common need: precise theoretical predictions. In particle physics, differential cross sections serve this purpose, providing the basis for testing the Standard Model and searching for possible hints of new physics. In gravitational-wave physics, the analogous role is played by the waveforms emitted by coalescing compact binaries, whose accurate modelling is essential both for the detection of signals and for the extraction of the astrophysical information they carry. From a computational point of view, these two subjects also share a significant amount of common structure.

The primary focus of this thesis is the study of the classical gravitational dynamics of two compact objects in different astrophysical situations, in particular the \textit{inspiral} and \textit{scattering} regimes, within effective theories of gravity beyond General Relativity. The central aim is to compute the classically relevant observables in these settings using modern techniques from quantum field theory. The starting point is the recognition that, despite its remarkable success in the weak-field regime, General Relativity is not expected to be the final or UV-complete description of gravity. As a classical theory, it should eventually arise from a consistent quantum framework, while as a low-energy theory, it is naturally understood as the leading term in a more general effective expansion. This viewpoint motivates the study of effective theories of gravity containing higher-curvature corrections and additional degrees of freedom. Such extensions typically introduce new propagating modes beyond the standard massless spin-2 graviton, and therefore lead to departures from the dynamics of pure General Relativity.

On the other hand, Quantum Field Theory (QFT) was originally developed to describe the interactions of fundamental particles. Over the last decade, however, it has become increasingly clear that many techniques borrowed from QFT are also extremely powerful for computing classical observables, provided the external states are prepared appropriately. In the standard particle-physics setting, perturbation theory is often described as an expansion both in the couplings and in $\hbar$. While this viewpoint is useful for scattering of microscopic quantum particles, it is not strictly correct when the external states correspond to macroscopic classical objects such as black holes or neutron stars. In such cases, classical information is not obtained simply by keeping tree-level terms and discarding loops; rather, it emerges from suitable kinematical limits of amplitudes and correlators, for instance through soft-momentum expansions in the large-impact-parameter regime.

There exist several complementary approaches for addressing these problems in different physical regimes. For the inspiral problem, where the orbital motion is non-relativistic, one employs the framework of Post-Newtonian Effective Field Theory (PN-EFT), also known as Non-Relativistic General Relativity (NRGR). The key idea is to exploit the hierarchy of length scales in the binary system, namely the size of the compact object, the orbital separation, and the wavelength of the emitted radiation. This separation of scales allows one to construct a systematic effective field theory description in which the conservative and radiative sectors can be treated in an organized manner.

In contrast, for scattering processes involving compact objects moving at relativistic velocities, the post-Newtonian expansion is no longer adequate. In this regime, it is more natural to use the post-Minkowskian (PM) expansion, which is organized as an expansion in Newton's constant $G$ while retaining the full velocity dependence. Several QFT-based methods have been developed for this purpose. These include direct amplitude-based approaches, such as the PM expansion of amplitudes and the KMOC formalism, as well as worldline-based frameworks such as PM-EFT and Worldline Quantum Field Theory (WQFT), which is closely related to the eikonal scattering amplitude. Together, these methods provide a modern and efficient toolkit for extracting conservative and radiative observables from classical gravitational dynamics beyond General Relativity.

As a further direction, this thesis also explores the celestial representation of scattering amplitudes, in which four-dimensional amplitudes in asymptotically flat spacetime are recast as conformal correlators on the celestial sphere at infinity. A well-known difficulty in this subject is that, in perturbative quantum field theory, especially for UV-incomplete theories, celestial amplitudes are often divergent order by order in perturbation theory. One possible way to improve this behavior is to resum the perturbative series before performing the Mellin transform, a strategy that becomes practically feasible in the eikonal regime. In the context of General Relativity, such an eikonal resummation is known to significantly improve the analytic properties of celestial amplitudes. When higher-derivative corrections to General Relativity are included, the ultraviolet behavior of the theory is modified, and this in turn can lead to an improved analytic structure for the corresponding celestial amplitudes. At the same time, however, constructing eikonal amplitudes beyond General Relativity becomes increasingly challenging. In this thesis, we also address these difficulties by employing techniques borrowed directly from analytic functional analysis.

